\NormalFont\ShortTitle{Colossenses}
{\MT Carta de Paulo aos Colossenses

\par }\ChapOne{1}{\PP \VerseOne{1}Paulo, apóstolo de Jesus Cristo pela vontade de Deus, e o irmão Timóteo,
\VS{2}aos santos e fiéis irmãos em Cristo, que estão em Colossos.
{\ADD{Haja}} em vós a graça e paz da parte de Deus nosso Pai, e do Senhor Jesus Cristo.
\VS{3}Agradecemos ao Deus e Pai do nosso Senhor Jesus Cristo, sempre orando por vós;
\VS{4}pois ouvimos falar da vossa fé em Cristo Jesus, e do amor
{\ADD{que tendes}} para com todos os santos;
\VS{5}por causa da esperança que está reservada para vós nos céus. Anteriormente ouvistes dessa esperança
\FTNT{*}{{\FR 1:5 }
a palavra “esperança” foi repetida para simplificar uma frase longa
} pala palavra da verdade do evangelho,
\VS{6}que já chegou a vós, como também em todo o mundo; e vai frutificando assim como entre vós, desde o dia em que ouvistes e conhecestes a graça de Deus em verdade.
\VS{7}Aprendestes o evangelho
\FTNT{†}{{\FR 1:7 }
a palavra “evangelho” foi repetida para simplificar uma frase longa
} por Epafras, nosso amado cooperador, que é um fiel servidor de Cristo em vosso favor.
\VS{8}Ele também nos declarou o vosso amor no Espírito.
\VS{9}Por isso também, desde o dia que
{\ADD{o}} ouvimos, não paramos de orar por vós, e de pedir que sejais cheios do conhecimento da sua vontade, em toda sabedoria e entendimento espiritual;
\VS{10}para que possais andar dignamente
{\ADD{diante}} do Senhor, agradando-lhe em tudo, frutificando em toda boa obra, e crescendo no conhecimento de Deus;
\VS{11}capacitados em todo fortalecimento, segundo o poder da sua glória, em toda perseverança e paciência com alegria,
\VS{12}agradecendo ao Pai, que nos capacitou a participar da herança dos santos na luz.
\VS{13}Ele nos tirou do domínio das trevas, e nos transportou para o reino do seu amado Filho,
\VS{14}em quem temos a libertação pelo seu sangue: o perdão dos pecados.
\VS{15}Ele é a imagem do Deus invisível, o primogênito de toda a criação;
\VS{16}porque nele foram criadas todas as coisas que há nos céus e na terra, visíveis e invisíveis, sejam tronos, sejam domínios, sejam governos, sejam autoridades; todas as coisas foram criadas por ele e para ele.
\VS{17}Ele existe antes de todas as coisas, e nele todas as coisas se mantêm.
\VS{18}E ele é a cabeça do corpo, da Igreja; ele é o princípio
{\ADD{e}} o primogênito dentre os mortos, para que tenha a primazia entre todos.
\VS{19}Pois foi do agrado
{\ADD{do Pai}} que toda a plenitude habitasse nele;
\VS{20}e, havendo por ele feito a paz pelo sangue da sua cruz, por meio dele reconciliasse consigo mesmo todas as coisas, tanto as que estão na terra, como as que estão nos céus.
\VS{21}Vós éreis separados
{\ADD{dele}} , e inimigos no entendimento, em más obras. Porém agora ele vos reconciliou,
\VS{22}no seu corpo de carne, pela morte, para vos apresentar diante dele santos, irrepreensíveis e inculpáveis;
\VS{23}se, de fato, permaneceis fundados e firmes na fé, e não vos afastais da esperança do evangelho que ouvistes, que é pregado a toda criatura que há abaixo do céu, e do qual eu, Paulo, fui feito servidor.
\VS{24}Agora me alegro nos meus sofrimentos por vós, e cumpro na minha carne o restante das aflições de Cristo, em benefício do seu corpo, que é a Igreja.
\VS{25}Dela eu fui feito servidor segundo a responsabilidade da parte de Deus que me foi dada para o vosso benefício, a fim de cumprir a palavra de Deus,
\VS{26}o mistério que estava oculto desde os princípios dos tempos e das gerações; mas agora foi manifesto aos seus santos.
\VS{27}A eles Deus quis fazer que conhecessem as riquezas da glória deste mistério entre os gentios, que é Cristo em vós, a esperança da glória.
\VS{28}Ele é quem anunciamos, advertindo a toda pessoa, e ensinando a toda pessoa, com toda sabedoria; a fim de que apresentemos toda pessoa
{\ADD{como}} completa em Cristo Jesus.
\VS{29}Para isso eu também trabalho, lutando conforme a eficácia dele, que age em mim com poder.
\FTNT{‡}{{\FR 1:29 }
eficácia (atuação, poder, habilidade) e age, no grego, são palavras semelhantes
}

\par }\Chap{2}{\PP \VerseOne{1}Pois quero que saibais como é grande a luta que tenho para o benefício de vós, dos que
{\ADD{estão}} em Laodiceia, e de tantos quantos não viram meu rosto fisicamente;
\FTNT{*}{{\FR 2:1 }
fisicamente
Lit. em carne
}
\VS{2}para que os seus corações sejam consolados, e estejam unidos em amor, e em todas as riquezas da pleno entendimento, para conhecimento do mistério de Deus, e do Pai, e de Cristo.
\VS{3}Nele estão ocultos todos os tesouros da sabedoria e do conhecimento.
\VS{4}Digo isso para que ninguém vos engane com palavras persuasivas.
\VS{5}Pois ainda que esteja ausente em corpo, todavia em espírito estou convosco, alegrando-me, e vendo a vossa ordem e a firmeza de vossa fé em Cristo.
\VS{6}Portanto, assim como recebestes Cristo Jesus, o Senhor, assim também andai nele;
\VS{7}com raízes firmes e edificados nele, e confirmados na fé, como fostes ensinados, abundando nela com gratidão.
\VS{8}Cuidado que ninguém vos tome cativos por meio de filosofias e de enganos vãos, conforme a tradição humana, conforme os elementos do mundo, e não conforme Cristo.
\VS{9}Pois nele habita corporalmente toda a plenitude da divindade.
\VS{10}E vós vos tornais plenos nele, que é o cabeça de todo governo e autoridade.
\VS{11}Nele também fostes circuncidados com uma circuncisão não feita por mãos,
{\ADD{mas sim}} , no abandono do corpo dos pecados da carne, pela circuncisão de Cristo;
\VS{12}sepultados com ele no batismo, no qual também com ele fostes ressuscitados pela fé no poder de Deus, que o ressuscitou dos mortos.
\VS{13}E quando vós estáveis mortos em pecados, na incircuncisão da vossa carne,
{\ADD{Deus}} vos deu vida juntamente com ele, perdoando-vos todas as ofensas.
\VS{14}Ele riscou a certidão de nossa dívida em ordenanças, a qual era contra nós, e a removeu, cravando-a na cruz;
\VS{15}ele despojou os domínios e autoridades, publicamente os envergonhou, e nela triunfou sobre eles.
\VS{16}Portanto, ninguém vos julgue pelo comer, ou pelo beber, ou por causa de celebração
{\ADD{religiosa}} , ou pela lua nova, ou pelos sábados;
\VS{17}essas coisas são a sombra das coisas futuras, mas a realidade
\FTNT{†}{{\FR 2:17 }
a realidade
Lit. “o corpo”, mas quando está em contraste de sombra, a palavra se refere ao que é concreto, isto é, à realidade.
} pertence a Cristo.
\VS{18}Ninguém, pois, se faça de juiz contra vós, insistindo nos pretextos de “humildade” e de “culto aos anjos”
\FTNT{‡}{{\FR 2:18 }
Ou: “culto dos anjos”
} , envolvendo-se em coisas que nunca viu, vangloriando-se em sua mentalidade carnal,
\VS{19}e não se ligando à cabeça, da qual todo o corpo, suprido e organizado pelas juntas e ligamentos, vai crescendo com o crescimento da parte de Deus.
\VS{20}Portanto, se estais mortos com Cristo para os elementos do mundo, por que vos sujeitais a ordenanças como se vivêsseis no mundo,
\VS{21}{\ADD{tais como}} : “não pegues”, “não proves”, “não toques”?
\VS{22}Todas elas se acabam com o uso, e se baseiam em mandamentos e doutrinas humanas.
\VS{23}Elas realmente têm aparência de sabedoria em culto voluntário, humildade, e tratamento severo ao corpo; mas não têm valor algum contra a satisfação da carne.

\par }\Chap{3}{\PP \VerseOne{1}Portanto, se fostes ressuscitados com Cristo, buscai as coisas de cima, onde Cristo está sentado à direita de Deus.
\VS{2}Pensai nas coisas de cima, não nas que
{\ADD{estão}} na terra.
\VS{3}Pois já morrestes, e a vossa vida está escondida com Cristo em Deus.
\VS{4}Quando Cristo,
{\ADD{que é}} a nossa vida, se manifestar, então também vós vos manifestareis com ele em glória.
\VS{5}Portanto, mortificai as partes carnais que há em vós que são terrenas;
{\ADD{são elas}} : o pecado sexual, a impureza, a paixão depravada, o desejo maligno, e a ganância, que é idolatria.
\VS{6}Por causa delas que a ira de Deus vem sobre os filhos da desobediência.
\VS{7}Nessas coisas também vós andastes antes, quando vivíeis nelas.
\VS{8}Mas agora, abandonai, vós também, todas
{\ADD{estas coisas}} : a ira, a fúria, a malícia, a maledicência, e as palavras vergonhosas da vossa boca.
\VS{9}Não mintais uns aos outros, pois já vos despistes do velho ser humano com os seus costumes,
\VS{10}e vos revestistes do novo, que se renova em conhecimento, conforme a imagem daquele que o criou.
\VS{11}Nisso não há grego, nem judeu, nem circunciso, nem incircunciso, bárbaro, cita, escravo,
{\ADD{nem}} livre; mas Cristo é tudo, e em todos.
\VS{12}Por isso, como escolhidos de Deus, santos e amados, revesti-vos de sentimentos de misericórdia, bondade, humildade, mansidão, e paciência.
\VS{13}Suportai-vos uns aos outros, e perdoai-vos uns aos outros, se alguém tiver queixa contra outro. Como Cristo vos perdoou, assim também
{\ADD{fazei}} .
\VS{14}E acima de tudo isto,
{\ADD{revesti-vos}} do amor, que é o vínculo da perfeição.
\VS{15}E a paz de Deus, para a qual também fostes chamados em um
{\ADD{só}} corpo, governe em vossos corações, e sede gratos.
\VS{16}A palavra de Cristo habite abundantemente em vós com toda sabedoria, ensinando-vos e advertindo-vos uns aos outros; e cantai ao Senhor com salmos, hinos, e cânticos espirituais, com graça no vosso coração.
\VS{17}E tudo quanto fizerdes, por palavras ou por obras,
{\ADD{fazei}} tudo no nome do Senhor Jesus, dando por meio dele graças ao Deus e Pai.
\VS{18}Mulheres, sede submissas aos vossos próprios maridos, como convém no Senhor.
\VS{19}Maridos, amai às esposas, e não as trateis amargamente.
\VS{20}Filhos, obedecei em tudo aos pais, porque isso é agradável ao Senhor.
\VS{21}Pais, não provoqueis aos vossos filhos, para que não percam o ânimo.
\VS{22}Escravos, obedecei em tudo aos
{\ADD{vossos}} senhores segundo a carne, não servindo apenas quando observados, como que para satisfazer às pessoas, mas sim, com sinceridade de coração, temendo a Deus.
\VS{23}E tudo quanto fizerdes, fazei de coração, como para o Senhor, e não para as pessoas.
\VS{24}Pois sabeis que recebereis do Senhor a recompensa da herança, porque servis ao Senhor Cristo.
\VS{25}Mas aquele que fizer injustiça receberá conforme a injustiça que cometer; e não há acepção de pessoas.

\par }\Chap{4}{\PP \VerseOne{1}Senhores, procedei de maneira justa e imparcial com os
{\ADD{vossos}} servos, sabendo que também vós tendes um Senhor nos céus.
\VS{2}Perseverai na oração, vigiando nela com gratidão.
\VS{3}Orai também ao mesmo tempo por nós, para que Deus nos abra uma porta para a palavra, para que falemos do mistério de Cristo,
{\ADD{mistério}} pelo qual estou preso.
\VS{4}Para que eu o anuncie claramente,
\FTNT{*}{{\FR 4:4 }
Ou: abertamente
} como devo falar.
\VS{5}Convivei com sabedoria com os que estão de fora,
\FTNT{†}{{\FR 4:5 }
Isto é, os que não são cristãos
} aproveitando o tempo oportuno.
\VS{6}A vossa palavra
{\ADD{seja}} sempre com graça,
\FTNT{‡}{{\FR 4:6 }
com graça
Isto é, de maneira agradável
} temperada com sal, para que saibais como deveis responder a cada um.
\VS{7}Tíquico, irmão amado, fiel trabalhador e colaborador no Senhor, vos fará saber tudo o que vem acontecendo comigo.
\VS{8}Foi para isso que eu o enviei até vós, para que saiba como estais, e console os vossos corações.
\VS{9}Juntamente com ele
{\ADD{está}} Onésimo, irmão fiel e amado, que é um de vós. Eles vos farão saber tudo o que está havendo por aqui.
\VS{10}Aristarco, que está preso comigo, vos saúda; e também Marcos, primo de Barnabé. Acerca dele, já recebestes ordens; se ele for até vós, recebei-o.
\VS{11}E também Jesus, chamado Justo. Dentre os circuncisos, esses
{\ADD{são}} os únicos cooperadores
{\ADD{meus}} no Reino de Deus, e têm sido consolação para mim.
\VS{12}Epafras, que é um de vós, servo de Cristo, vos saúda. Ele sempre luta por vós em orações, para que fiqueis íntegros e plenos em toda a vontade de Deus.
\VS{13}Pois dou testemunho dele de que ele tem grande zelo por vós, pelos que
{\ADD{estão}} em Laodiceia, e pelos que
{\ADD{estão}} em Hierápolis.
\VS{14}Lucas, o médico amado, vos saúda; e também Demas.
\VS{15}Saudai aos irmãos que
{\ADD{estão}} em Laodiceia, e também a Ninfas, e à igreja que está na casa dele.
\VS{16}E depois que esta carta for lida entre vós, fazei que ela também seja lida na igreja dos laodicenses, e que vós leiais também a de Laodiceia.
\VS{17}E dizei a Árquipo: Presta atenção ao trabalho que recebeste do Senhor, para que o cumpras.
\VS{18}{\ADD{Esta}} saudação
{\ADD{é}} da minha mão, de Paulo. Lembrai-vos das minhas prisões. A graça
{\ADD{esteja}} convosco. Amém. (Escrita de Roma aos colossenses, e enviada por Tíquico e Onésimo.)\par }